\documentclass[11pt,a4paper]{article}
\usepackage[margin=2.6cm]{geometry}
\usepackage{amsmath,amssymb}
\usepackage{booktabs}
\usepackage[colorlinks=true,linkcolor=blue,urlcolor=blue]{hyperref}
\usepackage{parskip}
\usepackage{xcolor}
\newcommand{\intuition}[1]{\begin{quote}\small\itshape #1\end{quote}}

\title{Predicting the S\&P 500 Official Opening Print\\
\large A complete description of the model, the news layer, and the trading approach}
\author{SPX Open Predictor project}
\date{July 15, 2026 --- working draft for review}

\begin{document}
\maketitle

\begin{abstract}
Every trading day, Polymarket runs a market called ``S\&P 500 Opens Up or Down''.
It pays out based on one specific number: the \emph{official} opening value of the
S\&P 500 index, as published by the Wall Street Journal. This document explains,
in plain language, how our system predicts that number: a statistical core model
built on overnight futures, a planned reconstruction of the opening auction
itself, and a news layer driven by a language model. Every parameter mentioned
is fitted from roughly two years of historical data; nothing is hand-tuned.
The document ends with open questions where suggestions are most useful.
\end{abstract}

\tableofcontents

\section{The bet and the number it resolves on}

The Polymarket contract asks a simple question: will the official S\&P 500 open
be higher or lower than yesterday's official close? Shares of the winning side
pay \$1; the losing side pays \$0. Prices before resolution behave like
probabilities: if ``Up'' trades at \$0.70, the crowd collectively believes there
is a 70\% chance the index opens higher.

The subtlety --- and the entire reason this project exists --- is \emph{which}
number decides the outcome. It is not the futures price, not an ETF price, not
the level shown on a TradingView chart. It is the first index value that
Standard \& Poor's computes at 9:30:00 New York time, echoed by the WSJ.

\section{Why the official open is a strange number}

The S\&P 500 is a weighted average of 500 stock prices. At 9:30:00 sharp, most
of those stocks have not actually traded yet: each stock opens with its own
auction on its own exchange, and those auctions finish at slightly different
moments. S\&P's calculation does not wait. It computes the first index tick
using \textbf{yesterday's closing price for every stock that has not opened
yet}.

\intuition{Picture a choir where each singer joins one by one. The first chord
you hear is mostly yesterday's recording, with only a few live voices mixed in.
Nasdaq-listed stocks (about 45--57\% of index weight, including Apple, Nvidia,
Microsoft) join at exactly 9:30:00; many NYSE names lag by seconds to minutes.}

Two measurable consequences follow. First, the official open only realizes
about 80\% of the move that futures markets predict --- the stale stocks drag it
toward yesterday's close. Second, on days when futures point only slightly up
or down, the official open can land on the \emph{opposite} side of what every
live chart shows. We call these \textbf{auction-quirk days}; they occur roughly
once a month, and they are the days the Polymarket crowd loses most
spectacularly (we measured six such days in five months --- on all six, the
crowd held the losing side at 94--99 cents).

\section{The core market model}

\subsection{The idea}

Overnight, S\&P 500 futures (ES) trade continuously and absorb almost all news.
Our core model reads the futures ``gap'' --- the percentage difference between
the futures price now and yesterday's 16:00 close --- and converts it into a
probability that the official open prints higher.

\subsection{The formula}

Let $g$ be the futures gap and $\tau$ the minutes remaining before 9:30. The
model predicts the official gap as
\[
\mu_m \;=\; a_0 + k \cdot g \;+\; c_{nq}(h)\,(q_{NQ}-g),
\qquad
P(\text{up}) = \Phi\!\left(\frac{\mu_m}{\sigma_m}\right),
\]
where $\Phi$ is the normal distribution function. The pieces:

\begin{itemize}
\item \textbf{Attenuation $k \approx 0.8$.} The official open realizes only
part of the futures move (the stale-choir effect). Fitted, not assumed ---
and the crowd's core error is trading as if $k=1$.
\item \textbf{$k$ depends on conditions.} In calm regimes $k\approx0.73$; in
volatile regimes $k\approx0.84$ (news carries through more fully). For small
gaps early in the night, $k\approx1$.
\item \textbf{The Nasdaq spread term.} Because the first tick is dominated by
Nasdaq names, the spread between Nasdaq (NQ) and S\&P (ES) futures carries
extra information early in the night.
\item \textbf{Uncertainty $\sigma_m$ shrinks as the open approaches:} about
$\pm0.34\%$ at midnight, $\pm0.15\%$ at 9:00, $\pm0.08\%$ at 9:29. It is
scaled up on volatile nights using recent realized volatility and the VIX1D
index, and multiplied by fitted event factors described below.
\end{itemize}

\subsection{Measured accuracy}

Direction accuracy on out-of-sample data (probability computed, then compared
with the realized official open):

\begin{center}
\begin{tabular}{lccccc}
\toprule
Prediction time & 00:00 & 04:00 & 07:00 & 09:00 & 09:29 \\
\midrule
Accuracy & 72.5\% & 77\% & 85\% & 90\% & $\sim$94\% \\
\bottomrule
\end{tabular}
\end{center}

Even with perfect market data at 9:30:00 the ceiling is 96.3\%, because the
remaining error is auction mechanics, invisible to any futures-based method.
Crucially, the model \emph{knows when it does not know}: dropping nights where
its own confidence is below 65\% (about 4 nights in 10 at midnight) raises the
accuracy of the remaining calls to $\sim$82\% at midnight and $\sim$94\% at 9:00.

\subsection{Where the model fails --- a complete taxonomy}

We classified every wrong committed call over 465 trading days. Three causes
cover all of them:
\begin{enumerate}
\item \textbf{No signal existed} (coin-flip nights; excluded by the confidence
filter).
\item \textbf{News arrived after the prediction} ($\sim$half of all failures;
flip windows cluster at the European open and above all in the 8:30--9:30
release hour).
\item \textbf{Auction quirks} (futures held one side all night; the official
print landed on the other). Unfixable with futures data --- empirically
confirmed, which motivates the auction layer.
\end{enumerate}

\section{Layer R: reconstructing the opening auction (in progress)}

Both exchanges publish, before 9:30, per-stock \emph{indicative} auction prices
--- the exchange's own running estimate of where each opening auction will
print. Nasdaq's feed (``NOII'') updates every second from 9:28; NYSE's
imbalance feed runs from 8:00.

Using a free historical sample day we measured how close those indicative
prices come to the actual opening prints: for Nasdaq mega-caps, within
\textbf{2--5 basis points} at 9:29; aggregated to index level the error is
roughly \textbf{1--2 index points} --- versus $\sim$6 points for the futures
model at the same moment. That is the quirk-day killer: a weighted
reconstruction of the first tick,
\[
\mu_R = \sum_i w_i \left(\frac{\hat{P}_i}{C_i} - 1\right)\cdot
\mathbf{1}[\text{stock } i \text{ in the first tick}],
\]
where $\hat{P}_i$ is the indicative price, $C_i$ yesterday's close, and stale
stocks contribute exactly zero. NYSE names default to ``stale'' --- we measured
their indicative prices to be 6--10$\times$ noisier than Nasdaq's.

\subsection{Implementation via LSEG Tick History}

LSEG Tick History (accessed through the DataScope Select REST API) is an
archive of every message the exchanges published, going back to 1996 ---
including the per-second Nasdaq NOII indicative prices and the NYSE imbalance
feed for any past date. Access has been approved through a university
(academic license); credentials are pending. The pipeline is already coded and
waits only for the login:

\begin{enumerate}
\item \textbf{Authenticate} (\texttt{python3 . lseg-status}) --- exchanges the
university credentials for an API token.
\item \textbf{Pull one morning} (\texttt{lseg-pull -{}-date 2026-06-22}) ---
requests all $\sim$505 index constituents, both exchanges, for the 9:15--9:40
window of that date; the server prepares an extract which we download and
cache ($\sim$tens of MB per day).
\item \textbf{Field discovery} --- LSEG stores each exchange field under its
own name; the first real pull prints a frequency table of the field names seen
so the parser's mapping can be confirmed once, by hand, against the spec.
\item \textbf{Replay} (\texttt{stage3 -{}-date 2026-06-22}) --- rebuilds the
first index tick from the indicative prices via the $\mu_R$ formula above and
compares it against the official open: one line of output per day, MATCH or
MISS, with the replica's error in index points.
\item \textbf{Batch validation} --- repeat over the 12 known quirk days plus a
control sample of normal days; fit the replica's real uncertainty
$\sigma_R(\tau)$ from these replays (currently a placeholder prior), which is
the number the fusion rule needs.
\end{enumerate}

One caveat belongs in print: the academic license covers research --- i.e.\
exactly this validation --- but a production system trading live must buy the
same data commercially (dxFeed quote, or NYSE TAQ at \$1{,}000 for a
12-months-back subscription).

\subsection{What this should theoretically deliver}

The free-sample measurement (indicative-vs-print error of 2--5\,bp per
mega-cap, $\sim$1--2 index points aggregated) lets us state the expected
outcome before spending anything:

\begin{center}
\begin{tabular}{lll}
\toprule
 & Today (futures only) & With validated replica \\
\midrule
9:29 direction accuracy & $\sim$94\% & $\sim$97--98\% (ceiling no longer binds) \\
Index-level error at 9:29 & $\pm$6 pts & $\pm$1--2 pts; $\pm$0.75 at 9:29:50 \\
Quirk days handled & 2 of 12 & majority called; rest honestly priced \\
Quirk-day Polymarket stance & stand down & buy the 1--6\textcent\ side the crowd abandoned \\
\bottomrule
\end{tabular}
\end{center}

The honest nuance: the very smallest quirk days (official gap of
$\pm$0.1--0.2 points, like June 22's $-0.14$) sit \emph{inside} even the
replica's error band --- those days the replica will not call with certainty
either. But it does not need to. The crowd prices such days at 94--99 cents;
the replica prices them at their true 60--80\%. Buying the abandoned side at
1--6 cents with a true probability near 30--40\% is a several-fold
expected-value trade \emph{even when the replica cannot name the winner}. On
the larger quirk days (0.5--0.9 points, most of the twelve), a $\pm$0.75-point
estimate at 9:29:50 is decisive. The 96.3\% ``ceiling'' of
Section~3.3 stops applying entirely, because it was the limit of
market-wide price data --- the replica reads the auction book itself, which is
the mechanism that generates those residual errors.

The second, quieter payoff is the fitted $\sigma_R$: once the replica's
uncertainty is measured on real replays, the fusion rule automatically hands
authority to it in the final two minutes --- no thresholds to tune --- and the
same code path runs live later with a commercial feed (NCDS/dxFeed for Nasdaq,
Massive for NYSE) without modification.

\subsection{What exactly is unknown at 9:30:00 --- the 505 stocks, precisely}

A common misreading is that ``which stocks are in the index'' varies. It does
not: \textbf{the membership and the weights are published and fixed} --- the
index is always the weighted sum of all $\sim$505 names. What varies, second
by second, is \emph{which price each stock contributes} to the first
calculation:

\begin{itemize}
\item a stock whose opening auction has already printed contributes its
\textbf{fresh opening price} (a ``live'' voice);
\item a stock whose auction has not finished contributes
\textbf{yesterday's close} (a ``stale'' voice).
\end{itemize}

The Nasdaq half ($\sim$250 names, 45--57\% of weight) is nearly
deterministic: the opening cross fires at exactly 9:30:00, so those names are
essentially always live. \textbf{The NYSE half is the stochastic part}:
DMM-run auctions straggle by seconds to minutes, and which NYSE names make it
into the first tick changes day by day and stock by stock. Our fitted
attenuation $k\approx0.78$ implies that on an average day roughly 20--25\% of
index weight \emph{beyond} the Nasdaq names does sneak into the first tick
--- the fast-opening NYSE mega-caps --- which our current replica models with
the blunt rule ``NYSE = always stale''. On the tiniest quirk days (decided by
$\pm$0.1--0.2 index points), whether three fast NYSE giants like JPM or LLY
are counted live can be the whole answer.

\subsection{Planned machine-learning upgrade of the replica}

The project has deliberately avoided machine learning so far: the day-level
problem offers only $\sim$500 observations (one open per day), where a
flexible model would memorise regimes instead of learning structure. LSEG
data breaks that constraint --- not by adding sources, but by changing the
observation unit. Per-second auction messages make the sample a
\emph{(stock, second)} pair: $505 \times \sim 120$ pre-open seconds
$\times$ 30 days $\approx$ \textbf{2 million samples}, which is genuine ML
territory. Two models, both trained from the same LSEG pulls:

\paragraph{Model 1: auction-residual model.} Predicts, per stock and per
morning, how far the final print will land from the current indicative price
--- replacing the single default correction and uncertainty with a
\textbf{dedicated per-stock value}. Features: near--far spread,
imbalance-to-paired ratio, price-variation indicator, seconds to the cross,
indicative-price velocity, and the stock's own historical auction noise.
Target: $\text{print} - \text{indicative}$ in basis points. Output: a bias
correction and a per-stock $\sigma_i$, which tighten the index-level
$\sigma_R$ and hence the replica's authority in the fusion.

\paragraph{Model 2: NYSE fast-open classifier.} Predicts, for each NYSE
stock each morning, the probability that its opening print exists in time to
be counted in the first tick. Features from the imbalance feed (size and
stability of the imbalance since 8:00, indicative-price volatility, the
stock's open-delay history, that morning's gap size); labels come free from
the print timestamps in the same data. The output enters the assembly not as
a hard switch but as an \emph{expected contribution}:
\[
\text{contribution}_i \;=\; w_i \Big[\, p_{\text{live},i}\cdot g_i \;+\;
(1-p_{\text{live},i})\cdot 0 \,\Big],
\]
so a stock that is 70\% likely to be counted contributes 70\% of its
predicted gap. This replaces the current constants ($p_{\text{live}}=1$ for
Nasdaq, $0$ for NYSE) with calibrated per-stock, per-morning probabilities.

\paragraph{Three guardrails.} (i) \emph{No end-to-end model}: a learner
trained on $\sim$500 day-level labels (``all data in, open direction out'')
would overfit and destroy the calibration the fusion depends on; the
day-level core stays linear. (ii) \emph{Tooling quarantine}: ML libraries
live in a research directory that produces fitted constants; the runtime
stays stdlib-only. (iii) \emph{License inheritance}: models trained on
academically-licensed data answer the research question (does the ML-boosted
replica beat the plain one, in basis points?); the production versions are
retrained on commercially-licensed data before live use.

\section{The news layer}

\subsection{The key empirical fact}

We verified on our own event days that \textbf{news publication and market
reaction are simultaneous} (within one 15-minute indexing cycle). Nobody reads
tomorrow's news early: when Trump posted about halting strikes on Iran, the
first article and the 2\% futures jump landed in the same hour. Two design
rules follow:

\begin{enumerate}
\item The language model can never front-run headlines; its value is
\emph{interpreting} events better than the crowd repositions, and knowing when
risk is elevated.
\item The LLM must never see market data (prices, gaps, our position). That
evidence is already in the core model; repeating it to the LLM would
double-count it and anchor the LLM's judgment. Market moves only \emph{trigger}
LLM runs; they are never LLM \emph{input}.
\end{enumerate}

\subsection{Group A: the midnight briefing}

Slow-moving context is fetched once, just before the first model run
($\sim$00:05): tomorrow's scheduled 8:30 releases, tonight's mega-cap earnings,
inflation nowcasts versus consensus, and baseline odds from large prediction
markets (Fed decisions, war-risk markets). We measured that these structural
odds move only $\sim$0.2 points overnight, so one snapshot stays fresh until
the open.

A strictly-formatted LLM run turns this briefing into two numbers: a variance
multiplier $m_A \in [1, 2.5]$ (anchored to a \emph{fitted} value of 1.29 for
release mornings --- the data almost exactly confirmed our 1.3 guess) and a
small directional tilt $|d_A| \le 0.3$, permitted only when a nowcast clearly
disagrees with consensus.

\intuition{Group A mostly answers ``how wild could tonight be?'', not ``which
way?''. Before a scheduled release, risk is symmetric --- the honest
contribution is width, not direction.}

\subsection{Group B: the continuous watch}

Unpredictable news needs continuous listening: a streaming financial wire
(sub-minute), polling of key political accounts (no official API exists; 2--5
minute scrapes), a global news sweep at each checkpoint, SEC filings in the
evening, and the official release endpoints at exactly 8:30:00. The LLM runs at
six scheduled checkpoints (02:30, 04:30, 07:00, 08:00, 08:35, 09:15 --- chosen
from the measured distribution of when decisive news actually happens) plus
event-triggered runs, capped at six per night.

Its fitted mathematical effects:
\begin{itemize}
\item \textbf{Post-shock widening.} After a gap move of $\ge0.35\%$ within one
hour, the official open is measurably harder to predict: residuals are
$\sim$1.9$\times$ wider (fitted). $\sigma_m$ is widened accordingly for two
hours.
\item \textbf{Conflicted-day direction --- the most valuable find.} When a
shock happened overnight but the gap has drifted back to nearly flat, the
\emph{direction of the shock} still predicts the official open with
\textbf{85\% accuracy} (81\% at 7:00, 93\% at 8:00; $n=41$). These flat-gap
mornings are exactly the days the core model calls a coin flip and the
Polymarket crowd misprices hardest. Only here does the LLM's direction enter:
\[
\mu_B = 0.9 \cdot d_B \cdot c \cdot \sigma_m
\qquad (d_B \in [-1,1] \text{ direction},\; c \in [0,1] \text{ confidence}),
\]
and on already-repriced days the directional term is \emph{zero by
construction} --- the gap has the news in it, and adding it again would count
the same evidence twice.
\end{itemize}

\section{Putting it together: the fusion rule}

Every layer produces an estimate with an uncertainty: the market core
$(\mu_m, \sigma_m)$, the news voices $(\mu_A, \sigma_A)$, $(\mu_B, \sigma_B)$,
and --- in the final two minutes, once live --- the auction replica
$(\mu_R, \sigma_R)$. They are combined by inverse-variance weighting (the
statistically optimal average, where more certain voices count more):
\[
\mu_{\text{fused}} = \frac{\sum_i \mu_i/\sigma_i^2}{\sum_i 1/\sigma_i^2},
\qquad
\sigma_{\text{fused}} = \Big(\sum_i 1/\sigma_i^2\Big)^{-1/2},
\qquad
P(\text{up}) = \Phi\!\left(\frac{\mu_{\text{fused}}}{\sigma_{\text{fused}}}\right).
\]

\intuition{No hand-offs, no if-then switches. At midnight the futures voice
dominates because it is the sharpest. If the auction replica comes online at
9:28 with a tenth of the uncertainty, it automatically takes over. News voices
have a floor on their uncertainty, so they can nudge and widen but never
outvote a confident market estimate.}

Trading rules on top: stand down when fused confidence is below 65\%; trade
against the Polymarket price only when our probability differs from it by at
least 5 cents.

\section{How the trading strategy was chosen: an execution-realistic backtest}

The layers above produce a probability. Turning that probability into orders
--- when to enter, at what price, at what size, when to bail out --- was
decided empirically, on the full history of the Polymarket market itself,
under an execution model calibrated from real trade prints. This section
documents that derivation: what was tested, what died, what survived, and the
analysis behind each choice. Nothing here is hand-tuned on the evaluation
data: every parameter was selected on the first (chronological) half of the
history and validated once on the untouched second half.

\subsection{The laboratory}

The daily market has existed since January 22, 2026. We fetched all 120
resolved markets with their full minute-level price curves and joined 115 of
them to the model's hourly out-of-sample probabilities (the model is fitted
only on pre-2026 futures data, so \emph{every} strategy day is out of
sample for the model). Chronological split: 57 train days (through April 15),
58 test days (April 16 -- July 14). Polymarket's taker fee
($0.04\,p(1-p)$ per share) began March 30, so the entire test half is
fee-era --- a built-in regime stress.

Execution realism was calibrated, not assumed. From $\sim$3{,}400 real trade
prints sampled across six days: effective half-spread $\approx$ 1 cent,
price impact $\approx$ 0.05 cents per \$100 of stake --- mild, because the
binding constraint turns out to be \emph{depth}, not slope (below). Maker
fills are modelled pessimistically (the resting price must trade
\emph{through} the limit), and an independent audit of the harness found and
fixed four critical defects before any conclusion was drawn: fictional fills
on zero-volume days (the price curve is book midpoints, which exist even
when nobody traded --- one such ``fill'' was 46\% of an early result's
profit), an entry gate that accidentally admitted bets \emph{against} the
model's own view, over-optimistic maker fills, and a look-ahead in the news
overlay's fitted constants (since refit on disjoint data). The lesson is
general: \textbf{in a thin market, the backtest's execution assumptions are
as load-bearing as the signal.}

\subsection{The graveyard: what 15 verified candidates taught us}

Six strategy families were designed and swept by independent optimizers on
the train half, then walk-forwarded. The result is brutal and informative:

\begin{center}
\begin{tabular}{lccl}
\toprule
Family & Train (\$100 start) & Test (walk-forward) & Verdict \\
\midrule
Ride crowd, flip late & profitable & \$100 $\to$ \$36 & dead \\
Cheap-side longshots & \$100 $\to$ \$129 & near ruin (\$5--7) & dead \\
Sell-into-convergence & \$100 $\to$ \$200 & \$100 $\to$ \$96 & dead \\
Fixed entry, hour $\ge 5$ & strong & negative & window closed by fees \\
Maker-limit entries & strong & unverifiable & midpoint fictions \\
Fixed entry at hour 4 & \$100 $\to$ \$155 & \$100 $\to$ \$136 & \textbf{sole survivor} \\
\bottomrule
\end{tabular}
\end{center}

\intuition{Everything that monetized the January--March regime --- huge
volumes, fee-free, lazy repricing --- evaporated out of sample. The only
edge that survived the market's own evolution is the simplest one: the model
knows the direction hours before the crowd finishes pricing it.}

\subsection{Entry trigger: the first hour confidence arrives}

The surviving fixed-hour entry leaves information on the table: days that are
coin flips at 4:00 but \emph{clarify later} (a release resolves, the gap
firms up). Replacing ``enter at hour 4'' with ``enter at the \emph{first}
hour in 4--9 at which the fused probability leaves the coin-flip zone''
(edge $\ge$ 5 cents \emph{and} traded-side probability $\ge$ 0.70) added
eight late entries on the test half --- all eight won (+\$108 per \$50
staked) --- lifting the walk-forward from \$136 to \$145 per \$100 flat, and
to \$322 per \$100 under compounding. Two warnings came with it: at a 0.60
gate the same idea goes to \emph{ruin} (days that clarify weakly are traps
--- fake post-release certainty at already-converged prices), and late
entries at 86--95 cents need a $\sim$90\% win rate just to break even, so
the gate is load-bearing, not cosmetic.

\subsection{Entry timing: why ``buy immediately'' beats hunting the dip}

The obvious objection --- surely one should wait for a better price after
the signal --- was tested three ways and closed with an attribution
analysis. For every train-half signal day we located the post-signal price
minimum of our side and asked \emph{why the market was lowest there}:

\begin{itemize}
\item On \textbf{winning days} the median post-signal dip is \textbf{2
cents}. The deep dips (up to 47 cents) are, in 12 of 26 cases,
\emph{futures-adverse}: ES genuinely moved against the position mid-morning
--- at the moment of the dip our own signal had degraded below the gate, so
the ``discount'' was correctly-priced worse odds, not a bargain.
\item On \textbf{losing days} the median ``dip'' is \textbf{62 cents}: the
minimum is simply the crash to zero. A resting buy order below the market is
\emph{guaranteed} to fill on these days.
\end{itemize}

That asymmetry is adverse selection in its purest form, and it shows up in
every rule we derived from it: dip-triggered taker entries, resting maker
limits at $-1/-2/-3$ cents with timed fallbacks, and 50/50 splits all
underperform immediate entry on \emph{both} halves (train +\$257 immediate
vs +\$184 best variant; test +\$214 vs +\$168). The conclusion is
structural, not parametric: \textbf{after a valid signal the market
converges away from you; the only reliably available discounts are moments
when the signal is invalid, or days when it is wrong.} Entry is therefore
immediate, as taker --- the 1-cent spread and the fee are the price of
being early, and being early is the entire edge.

\subsection{The price gate}

Immediacy exposed one genuine flaw: at very high prices the trade shape
degenerates (risk \$400 to win \$18 at 95.4 cents --- one quirk-day loss at
such prices costs 22 wins to repay). A price cap on entries was swept on the
train half: rejecting entries above 95 cents removes exactly these
last-cent trades at zero cost, while a 90-cent cap overshoots (it deletes
profitable low-90s days, $-\$27$ on train). The gate is asymmetric by
design: expensive entries have capped reward and uncapped (100\%) loss;
cheap entries have the opposite shape, which is why the same 0.70
confidence gate that protects 90-cent entries is allowed to pass 50-cent
ones.

\subsection{Sizing: the half-rule, the cap, and the market's own ceiling}

Stake economics were measured by sweeping flat stakes from \$25 to \$2{,}000
with the calibrated execution model. Three facts emerged. First, return per
dollar is \emph{flat} ($\sim$15.5\%) from \$25 to \$300: no soft bottleneck.
Second, above $\sim$\$400 the simulation leaves its validated region (the
impact model was calibrated on \$100--500 prints). Third --- the real
ceiling --- stakes are capped by \emph{depth}: at most $\sim$1\% of a day's
traded volume can be taken without fiction, and intended stakes asymptote to
an effective $\sim$\$1{,}000 regardless of intent. The sizing rule is
therefore
\[
\text{stake} \;=\; \min\bigl(\tfrac{1}{2}\,\text{balance},\; \$400,\;
1\%\times\text{day volume}\bigr),
\]
with sub-\$1k-volume days skipped entirely. Walk-forward from \$100: test
half \$851 (\$669 with all execution costs doubled), peak \$1{,}308,
maximum drawdown \$477. The half-rule's danger is stated plainly: each
resolution loss halves the balance. The same rule on the \emph{train} half
(five losses rather than two) turns \$100 into only \$212 --- the outcome
is hostage to loss frequency, and a losses-first ordering of the test trades
survives only at \$24 at the trough. Halving the fraction roughly halves
both the growth and the violence; the choice is risk appetite, not
mathematics.

\subsection{Position monitors: sell the corpse, sometimes flip}

Because resolution losses are total, the position is re-examined when new
information arrives: at the 08:35 checkpoint (after the 8:30 release) and
the 9:00 model refresh. If the traded side's fused probability has fallen
below 0.5, the position is sold at market; if the \emph{other} side then
clears the full entry test, the proceeds flip into it. Quantified on our two
test losses: the July 2 news reversal was recoverable at 8:29 for \$77 of
the \$296 lost (and nothing later --- the market repriced within minutes);
the July 10 loss was fully invertible (sell + flip at 8:29 yields
\emph{+\$347} against $-\$307$ held) --- but only with information the
hourly futures did not have. Which leads to the final two layers.

\subsection{The last-minute futures leg}

Minute-level futures history only survives $\sim$30 days back, so this leg
was validated on the 13 available days (June 25 -- July 14): the 9:29 gap,
scored as $\Phi(0.79\,\text{gap}/0.083)$, called the official open
\textbf{13 of 13} times. Two tradeable patterns appeared immediately: days
where the crowd lags a late futures move (June 30: market at 0.35 while
9:29 futures implied 0.80 --- worth +\$88 per \$50), and --- critically ---
July 10 itself, previously classified as an unknowable quirk, turned out to
be a \emph{late futures flip} (gap $-0.03$ at 9:00 $\to$ $+0.08$ at 9:29)
that a 9:28 check would have converted from $-\$307$ into a win via the
sell-and-flip. A live 9:28--9:29 leg costs nothing (one live quote) and is
now part of the strategy; it cannot be backtested further for lack of
historical minute data, so it runs paper-first. True auction quirks (June 1,
where the winning side traded at 1.8 cents at 9:29; June 22 at 3 cents)
remain invisible to futures at \emph{any} minute --- capturing those is
precisely Layer~R's job, and the measured ceiling there (+\$1{,}734 on a
single \$50 stake on June 1) dwarfs everything else in this section.

\subsection{The resulting strategy, in one place}

\begin{enumerate}
\item \textbf{00:05} Group-A briefing $\to$ $(m_A, d_A)$; first fusion.
\item \textbf{02:30--09:15} Re-fuse at each checkpoint and on triggers
(gap/odds/wire); Group-B voice per Section 5.
\item \textbf{Entry} at the first hour in 4--9 with edge $\ge$ 5c,
traded-side $P \ge 0.70$, and price $\le$ 95c: immediate taker buy.
\item \textbf{Stake} $= \min(\frac{1}{2}\text{balance}, \$400, 1\%\times
\text{volume})$; skip days below \$1k volume.
\item \textbf{Monitors} at 08:35 and 9:00: sell below $P=0.5$ on the traded
side; flip when the other side fully qualifies.
\item \textbf{9:28} live-futures re-check: enter edges $\ge$ 3c; sell-and-flip
positions the minute-gap contradicts. (Layer~R replaces this voice when live.)
\item \textbf{9:30} resolution; log fills for execution recalibration;
monthly walk-forward refit.
\end{enumerate}

Validated expectation (58 out-of-sample fee-era days): \$100 $\to$ \$851
base / \$669 stressed; 26W--2L at 93\% win rate; both losses understood
(one news reversal, partially recoverable; one late futures flip, fully
recoverable with the 9:28 leg). Threats to validity, stated honestly: one
market regime, 58 evaluation days, midpoint price curves (maker-side
untrustworthy), a declining-volume market, and a crowd that can learn.

\section{Current status}

\begin{center}
\begin{tabular}{llp{6.2cm}}
\toprule
Layer & Status & Missing \\
\midrule
Market core & live, validated & --- \\
Confidence filter \& edge rule & live & --- \\
Auction replica code & built, sample-validated & LSEG credentials (approved, pending) for quirk-day replay; live feed choice \\
Group A prefetch & live (odds, calendar) & FRED key (free) for CPI/PPI dates; nowcast fetcher \\
Group A/B prompts \& math & written, fitted & LLM API runner \\
Group B feeds & designed & Alpaca key (free); Truth Social scraper choice \\
Replay evaluation (212 news days) & designed & the two keys above \\
\bottomrule
\end{tabular}
\end{center}

\section{Open questions --- where suggestions help most}

\begin{enumerate}
\item \textbf{Sizing.} We currently plan flat stakes above the confidence
threshold. Kelly-style sizing scaled by (fused confidence $-$ market price) is
the natural next step --- how aggressive is acceptable?
\item \textbf{The post-shock widening factor} (1.9) is unstable between sample
halves (1.3 vs 2.3). More data or a smarter conditioning variable?
\item \textbf{Conflicted-day sample size} is $n=41$ at 85\%. Enough to trade
at reduced size, or wait for live confirmation?
\item \textbf{Quirk-day product.} On tiny-gap days, is the better trade our
model's honest probability versus the crowd's 94--99 cent overconfidence
(buying the cheap side), or should we wait for the auction replica before
touching those days at all?
\item \textbf{LLM choice and cost.} $\sim$12 runs/night with strict JSON
outputs; latency matters only for the 08:35 run. Small fast model with strong
guardrails vs larger model at checkpoints only?
\item \textbf{Volume decline.} The Polymarket market's daily volume has fallen
from \$600k+ (January) to \$30--45k. The edge math still works but capacity is
shrinking --- worth monitoring whether this strategy stays worth the
infrastructure.
\end{enumerate}

\appendix
\section{Fitted constants and provenance}

\begin{center}
\small
\begin{tabular}{llll}
\toprule
Constant & Value & Meaning & Fitted from \\
\midrule
$a_0$ & 0.0092 & regression intercept (\%) & 6 months daily \\
$k$ (pooled) & $\approx$0.77--0.80 & futures$\to$official attenuation & 465d, train half \\
$k$ by regime & 0.73 / 0.83 / 0.84 & calm / mid / volatile tercile & 465d, train half \\
$\sigma$ at midnight & $\approx$0.34\% & residual width & hourly ES, 200d+ \\
$m_{\text{cal}}$ & 1.29 & release-morning widening & 11 NFP mornings (train) \\
$m_{\text{shock}}$ & 1.9 (1.3--2.3) & post-shock widening & 76 shock days \\
$Z_{\text{conflict}}$ & 0.9 & conflicted-day direction weight & 41 days, 85.4\% hit \\
NOII accuracy & 2--5 bp & indicative vs print, mega-caps & free ITCH sample \\
Ceiling & 96.3\% & perfect-market-data accuracy & 81 unseen days \\
\bottomrule
\end{tabular}
\end{center}

\section{Reproducing every number}

All results regenerate from the repository (\texttt{python3 .\ <command>}):
\texttt{backtest} (model metrics, quirk table), \texttt{news-mistakes}
(465-day news inventory), \texttt{news-timing} (decisive-hour profile),
\texttt{news-groupb-fit} (shock/conflict fits), \texttt{news-prefetch}
(tonight's Group-A block), \texttt{noii-deviation} (auction accuracy from a
raw ITCH file), \texttt{ground-truth} (official-open source audit).

\end{document}
